\documentclass[a4paper]{article}     
\usepackage{amsfonts}
\usepackage{amsmath}
\usepackage{amssymb}
\usepackage{graphicx}
\usepackage{color}

\newcommand{\dint}{\displaystyle\int}

\begin{document}

\noindent \large Auteur: Thomas Olszowka \\
\noindent \large Studierichting: Informatica \\
\noindent \large Oefeningenreeks 4A nr. 14.2  \\



\medskip

\normalsize

$I = \dint x \ln x dx$ \\

$\left\{\begin{array}{cc}
    u =& \ln x \\ 
    dv =& xdx  \\
\end{array}\right. \Rightarrow\left\{\begin{array}{cc}
    du =& \dfrac{1}{x} dx \\
    v =& \dfrac{x^2}{2}  \\
\end{array}\right.$ \\

$I = \ln x \cdot \dfrac{x^2}{2} - \dint \dfrac{x^2}{2} \cdot \dfrac{1}{x} dx$ \\

$= \dfrac{1}{2}x^2 \ln x - \dint \dfrac{x^2}{2x} dx$ \\

$= \dfrac{1}{2}x^2 \ln x - \dfrac{1}{2} \dint x dx$ \\

$= \dfrac{1}{2}x^2 \ln x - \dfrac{x^2}{4} + c$ \\

\begin{thebibliography}{999}
\end{thebibliography}
\end{document} 